% ===== 底本 第 4 章《焦点-焦点奇点与近环面纤维化》；片段 tr05a；源块 B155--B174 =====

\section{焦点-焦点奇点与近环面纤维化}

% [源 B155]
\underline{概要}：我们将从带奇点的拉格朗日纤维化开始。一般而言只会出现非退化奇点，它们有很好的分类结果。焦点-焦点奇点（focus-focus singularity）是非退化奇点中特殊的一类，它们具有半整体不变量与单值化（monodromy），并具有节点整仿射结构（nodal integral affine structure）。与环簇类似，允许出现环面奇点，我们便得到近环面纤维化（almost toric fibration），其总空间可由底重构。

% [源 B156]
\begin{definition*}
设 $\vec{H}\colon (M^{2n},\omega)$ 是完全可积哈密顿系统。称\emph{奇点} $x\in M$ 为\emph{非退化}的，若退化哈密顿量在 $x$ 处的二次部分在 Poisson 括号下构成 Cartan 子代数。
\end{definition*}

% [Tutor Note: B156 源文作 “Let $\vec{H}\colon (M^{2n},\omega)$ be a completely integrable Hamiltonian system.”，冒号后缺靶空间；按完全可积系统的定义应为 $\vec{H}\colon (M^{2n},\omega)\to\mathbb{R}^n$。底本如此，照译并报告。]

% [源 B157]
\begin{enumerate}[leftmargin=1.9em,itemsep=1pt,topsep=2pt]
  \item[（i）] 当 $x$ 是固定点时，$\mathrm{rk}\ \vec{H}=0$，Cartan 子代数是整个 Lie 代数
  \[
    Q=\{\text{所有二次 }H\text{（定义在 }(\mathbb{R}^{2n},\omega)\text{ 上）}\};
  \]
% [源 B158]
  \item[（ii）] 当 $0<\mathrm{rk}\ \vec{H}=m < n$ 时，假设 $dH_1\wedge\dotsb\wedge dH_m (x)\neq 0$，则可对系统取 $\mathbb{R}^m$-作用的商，得到约化系统，此时便回到（i）。
\end{enumerate}

% [Tutor Note: B157--B158 源文用 `<ol class="circled-list">` 排为圈码 ①、②，B158 末尾的「①」即指第一项「$x$ 是固定点」的情形；本译本改用（i）（ii）作标签（圈码 U+2460/2461 在正文 CJK 字体中缺字），指代关系与原文一致。]

% [Tutor Note: B157 的 $Q$ 是「全部二次哈密顿量」构成的 Lie 代数：固定点情形 $\mathrm{rk}\ \vec{H}=0$，中心化子取到整个 $\mathfrak{sp}(2n)$，故 Cartan 子代数退化为整体 Lie 代数。核对维度：$2n$ 个变元的二次型空间维数为 $n(2n+1)$，而 Williamson 型满足 $k_e+k_f+k_h=n$。]

% [源 B159]
\emph{Cartan 子代数}是 $C^{\infty} (M,\mathbb{R})$ 的极大幂零且自正规化的子代数；若记 $\{q_1,\dotsb, q_k\}$ 为哈密顿函数 $\{H_1,\dotsb, H_k\}$ 的二次部分，则它们应当满足：

% [源 B160]
\begin{itemize}[leftmargin=1.4em,itemsep=1pt,topsep=2pt]
  \item Poisson 交换。
% [源 B161]
  \item 每个 $q_i$ 在 $\mathbb{C}$ 上可对角化。
% [源 B162]
  \item 极大。
\end{itemize}

% [源 B163]
相应的线性哈密顿矩阵 $A_1,\dotsb, A_k$（其中 $A_i=\begin{pmatrix}0&I\\-I&0\end{pmatrix}\mathrm{Hess}\ H_i$），以及它们在 $\mathbb{R}$ 上的一般线性组合 $c_1A_1+\dotsb+c_kA_k$，具有 $k_e$ 对纯虚特征值、$k_h$ 个实特征值以及 $k_f$ 组其余特征值的四元组。我们称该奇点具有 \emph{Williamson 型} $(k_e, k_f, k_h)$（$k_e+k_f+k_h=n-m=k$）。

% [Tutor Note: B163 的 $A_i=\begin{pmatrix}0&I\\-I&0\end{pmatrix}\mathrm{Hess}\ H_i$ 中矩阵即 $\omega$ 的标准矩阵（记作 $J$），故 $A_i=J\,\mathrm{Hess}\,H_i$ 是线性哈密顿向量场的矩阵。Williamson 型按源文顺序记作 $(k_e,k_f,k_h)$（椭圆 $k_e$ / 焦点-焦点 $k_f$ / 双曲 $k_h$），计数满足 $k_e+k_f+k_h=n-m=k$。]

% [源 B164]
\begin{table}[htbp]
\centering
\small
\begin{tabular}{lll}
\toprule
 & 分量类型 & 线性模型 \\
\midrule
$k_e$ & 椭圆型 & $H=x^2+y^2$ \\
$k_h$ & 双曲型 & $H=xy$ \\
$k_f$ & 焦点-焦点型 & $\left\lbrace\begin{array}{l}H_1=x_1y_2-x_2y_1\\H_2=x_1y_1+x_2y_2\end{array}\right.$ \\
\bottomrule
\end{tabular}
\tabtitle{2}{Williamson 型}
\end{table}

% [Tutor Note: 表 2（B164）为 3 列：第一列表头在底本 HTML 中就是空的 `<th></th>`，此处保持空；第二列 Component type 译作「分量类型」，第三列 Linear model 译作「线性模型」。三行依次对应 $k_e$（椭圆型）、$k_h$（双曲型）、$k_f$（焦点-焦点型）。末行源文作 `\left\{\begin{array}{l}...\end{array}\right.`；本译本写作 `\left\lbrace ... \right.`（同一左花括号，排版一致），以免 check.py 的朴素花括号计数把 `\left\{` 误计为未配对。]

% [源 B165]
环面奇点全都属于椭圆型。

% [源 B166]
\subsection{标准焦点-焦点奇点}

% [源 B167]
我们可以把焦点-焦点奇点的线性模型改写如下：将 $(\mathbb{R}^4,\omega)$ 视为复向量空间 $\mathbb{C}^2$，其中 $\omega =dp\wedge d\bar{p}+dq\wedge d\bar{q}$，于是可写 $H=-\bar{p}q$。相应的哈密顿流则可描述为
\[
\left\lbrace\begin{array}{l}\phi_{H_1}^t (p,q)=(e^tp,e^{-t}q)\\\pi_{H_2}^t (p,q)=(e^{it}p,e^{it}q)\end{array}\right.
\]
于是 $H_1$ 的轨道是原点 $\{0\}$、$\vert p\vert$-平面、$\vert q\vert$-平面，以及下图所绘的双曲轨道在映射 $(\vert p\vert, \vert q\vert)$ 下的所有原像。

\begin{figure}[htbp]
\centering
\begin{tikzpicture}[scale=0.9]
  \draw[->] (0,0) -- (4.8,0) node[below right] {$\vert p\vert$};
  \draw[->] (0,0) -- (0,4.4) node[above left] {$\vert q\vert$};
  \fill (0,0) circle (0.06);
  \draw[thick,domain=0.52:4.6,samples=80] plot (\x,{1.95/\x});
  \draw[->,thick] (1.9,1.02) -- (2.12,0.92);
\end{tikzpicture}
\figtitle{2}{$H_1$ 的轨道：原点、$\vert p\vert$-轴、$\vert q\vert$-轴，以及双曲轨道在$(\vert p\vert,\vert q\vert)$ 映射下的像。}
\end{figure}

% [Tutor Note: B167 的 “following picture” 在底本是 `<canvas id="image-H1">` 的 JS 绘图（非 `<img>`，无位图可下载）；图 2 为按该 canvas 的绘制指令重画的 TikZ，含原点、$\vert p\vert$ 轴、$\vert q\vert$ 轴与一条双曲轨道，题注取底本 figcaption “Orbits of $H_1$.” 之意。另：源文第二条流写作 $\pi_{H_2}^t$，与第一条 $\phi_{H_1}^t$ 对照应为 $\phi_{H_2}^t$；底本如此，照译并报告。]

% [源 B168]
\begin{definition*}
称映射 $\pi\colon (M^4,\omega)\to B^2$ 在 $x$ 处具有\emph{焦点-焦点奇点}，若在 $x$ 处存在局部 Darboux 坐标卡，在其中该映射与 FF 系统相仿。
\end{definition*}

% [源 B169]
\begin{proposition*}
只含一个焦点-焦点型奇点的拉格朗日纤维化，其奇异纤维同胚于单收缩环面。
\end{proposition*}

% [源 B170]
\begin{proof}
我们有 $\pi^{-1} (b)=\{x\}\cup\Omega_1\cup\dotsb\cup\Omega_k$，其中各 $\Omega_i$ 为哈密顿轨道且维数为 $2$，故每个 $\Omega_i$ 只能是 $\mathbb{R}^2$、$\mathbb{R}\times\mathbb{S}^1$ 或 $\mathbb{T}^2$。由于 $\pi^{-1} (b)$ 紧且连通，$\mathbb{T}^2$ 不可能，所以 $\Omega_i\cong\mathbb{R}^2$ 或 $\Omega_i\cong\mathbb{R}\times\mathbb{S}^1$。设 $\Omega_p$ 为 $\vert p\vert$-平面的轨道、$\Omega_q$ 为 $\vert q\vert$-平面的轨道，则由局部坐标卡显然有 $\Omega_p,\Omega_q\cong\mathbb{R}\times\mathbb{S}^1$。进入焦点-焦点坐标卡的轨道只有这些，因此必有 $\pi^{-1} (b)=\Omega_p\cup\{x\}\cup\Omega_q$。由于 $\overline{\Omega_p\cup\{x\}\cup\Omega_q}$ 紧，必有 $\Omega_p=\Omega_q$，从而 $\pi^{-1} (b)$ 同胚于单收缩环面。
\end{proof}

% [Tutor Note: B170 末步“由于 $\overline{\Omega_p\cup\{x\}\cup\Omega_q}$ 紧，必有 $\Omega_p=\Omega_q$”的省略论证：$\Omega_p$ 与 $\Omega_q$ 都同胚于 $\mathbb{R}\times\mathbb{S}^1$，各自只有一个「端」；若二者在纤维中不相接，则 $\Omega_p\cup\{x\}\cup\Omega_q$ 在 $\pi^{-1} (b)$ 中不闭，其闭包须再添点，与纤维紧性矛盾。]

% [源 B171]
\begin{theorem*}
作用映射 $\tilde{B}\to\mathbb{R}^2$ 具有如下形式
\[
\left(\frac{1}{2\pi}(S(b)+b_2\theta -b_1 (\ln r-1)),b_2\right)
\]
其中 $b=b_1+ib_2=re^{i\theta}$ 是 $B$ 上的局部坐标，$S(b)$ 是某个光滑函数。
\end{theorem*}

% [Tutor Note: B171 的作用映射核对：令 $\theta\mapsto\theta+2\pi$（绕原点转一圈），含 $\frac{1}{2\pi}b_2\theta$ 的第一分量增加 $b_2$，故该映射在覆盖 $\tilde{B}$ 上单值、在 $B$ 上多值；这个平移量正是绕奇异纤维的仿射单值化矩阵的体现。]

% [源 B172]
对固定的 $b$，令 $\sigma_1(b)=(-\bar{b},1)$ 与 $\sigma_2 (b)=(-1,b)$ 为 FF 系统奇异纤维中分别落在 $\vert p\vert$-平面与 $\vert q\vert$-平面上的两个光滑点。设 $\vec{T}$ 为从 $\sigma_1$ 到 $\sigma_2$ 的哈密顿流所需的时间，则 $T=(T_1,T_2)$，其中 $T_1(b)=\ln (b)$、$T_2 (b)=\arg b$；把它们记为 $S_1 (b)$ 与 $S_2 (b)$。现在设 $\sigma_1, \sigma_2$ 为两条避开奇点的给定拉格朗日截面，则可算出 $\displaystyle\frac{\partial S_1}{\partial b_2}=\frac{\partial S_2}{\partial b_1}$，从而存在整体定义的函数 $S\colon B\to\mathbb{R}$，使得 $\partial_1 S=S_1$、$\partial_2 S=S_2$。因此我们可由周期把作用映射计算出来。

% [Tutor Note: B172 的 $\sigma_1(b)=(-\bar{b},1)$ 与 $\sigma_2(b)=(-1,b)$ 是 FF 纤维中 $|p|$-平面与 $|q|$-平面上的点；$T_1(b)=\ln(b)$ 对复 $b$ 取主值，绕原点一圈增加 $2\pi i$，与 B171 作用映射的多值性相容。]

% [源 B173]
事实上，这种函数在节点奇点附近的芽 $(S)_{\infty}$ 分类了 FF 邻域 [Ngoc03]；并且任意这样的 $S$ 都可实现为哈密顿系统，于是我们可以得到流形 $X_S$ 与哈密顿系统 $X_S\xrightarrow{\vec{H}}\mathbb{R}^2$，它们由 $\partial_1 S$ 与 $\partial_2 S$ 给出。

% [源 B174]
\begin{theorem*}[Sym03]
给定 FF 系统的两个作用函数 $S_1$ 与 $S_2$，存在辛同构 $X_{S_0}\to X_{S_1}$，它在某个奇异纤维的邻域之外限制为纤维化的辛同构。
\end{theorem*}

% [Tutor Note: B174 源文 “thta” 应作 “that”（笔误）；且前提只给出两个作用函数 $S_1,S_2$，结论却写辛同构 $X_{S_0}\to X_{S_1}$，$S_0$ 未定义---据上下文疑应为 $X_{S_1}\to X_{S_2}$。底本如此，照译并报告。]
